\documentclass[11pt]{article}
\usepackage[margin=1in]{geometry}
\usepackage{booktabs}
\usepackage{graphicx}
\usepackage{hyperref}
\usepackage{amsmath,amssymb}
\usepackage{natbib}
\usepackage{multirow}
\usepackage{caption}
\usepackage{subcaption}
\usepackage{xcolor}
\usepackage{algorithm}
\usepackage{algorithmic}

% Macro file generated by pipeline
\input{results.tex}

\title{Microstructure, Tail Risk, and Efficiency in Prediction Markets:\\
A Cross-Platform Empirical Study of 476 Million Trades}
\author{Andrey Orlov}
\date{\today}

\begin{document}
\maketitle

\begin{abstract}
Prediction markets have emerged as significant venues for price discovery on political, economic, and sporting events, yet the microstructure of these markets remains poorly understood. We provide the first large-scale empirical comparison of microstructure properties across a centralized, CFTC-regulated exchange (Kalshi) and a decentralized, on-chain market (Polymarket), analyzing over \NStandardizedTrades{} trades spanning a \NOverlapMonths-month overlap period. We address three research questions: (1) whether prediction market prices satisfy the martingale property after correcting for bid-ask bounce, (2) how tail risk characteristics differ by platform and market type under Extreme Value Theory, and (3) whether microstructure features can predict tail risk events with economic significance. Using variance ratio tests, Roll's implied spread model, Generalized Pareto Distribution fitting, BDS nonlinearity tests, and an automated model search over \NAutoresearchIters{} XGBoost configurations, we find that political markets are approximately weakly efficient after accounting for spreads, while sports markets harbor exploitable nonlinear structure (BDS $z$-scores of \BDSSportsZRange). Taker-side direction emerges as the dominant predictor of tail events (\TakerSideImportance\% feature importance), linking prediction market microstructure to the Probability of Informed Trading literature. A tail-risk-adjusted strategy achieves an out-of-sample Sharpe ratio of \BestSharpe{} on Kalshi, net of transaction costs.
\end{abstract}

% ============================================================================
\section{Introduction}
\label{sec:intro}
% ============================================================================

% [To be written by LLM discover/write loop. Key points:]
% - Prediction markets as information aggregation mechanisms
% - Two dominant platforms with fundamentally different designs
% - Gap: no large-scale microstructure study across centralized vs decentralized
% - Our contribution: 476M trades, three RQs, novel findings on efficiency/tails/predictability
% - Paper structure overview

% ============================================================================
\section{Related Work}
\label{sec:related}
% ============================================================================

% [Covers: prediction market efficiency (Tsangarides 2021, Arrow et al.),
%  market microstructure (Roll 1984, Lo-MacKinlay 1988, BDS 1996),
%  EVT in finance (McNeil-Frey 2000, Embrechts 1997, Coles 2001),
%  DEX microstructure (Adams et al. 2021, Uniswap),
%  prediction market empirics (Becker 2024, Ng et al. 2025, Reichenbach 2025, Arbutina 2025)]

% ============================================================================
\section{Data}
\label{sec:data}
% ============================================================================

Our dataset comprises \NStandardizedTrades{} standardized trade observations spanning two prediction market platforms.

\subsection{Kalshi (Centralized, Regulated)}

Kalshi is a CFTC-regulated event contract exchange operating a central limit order book. Binary contracts are priced in cents (0--99), with \texttt{yes\_price} + \texttt{no\_price} = 100 always. Our dataset contains \NKalshiTrades{} trades across \NKalshiMarkets{} unique markets, spanning June 2021 to November 2025.

\subsection{Polymarket (Decentralized, On-Chain)}

Polymarket operates on the Polygon blockchain using the Conditional Token Framework. Trade prices are derived from maker/taker amounts in USDC atomic units. Our dataset contains \NPolymarketTrades{} on-chain trades spanning March 2023 to January 2026. Timestamps are reconstructed by joining trades to block timestamps via \texttt{block\_number}.

\subsection{Data Processing}

% [Standardization: Kalshi price = yes_price/100, Polymarket price from amount ratios,
%  common schema: market_id, platform, price, quantity, side, timestamp]

% ============================================================================
\section{Methodology}
\label{sec:methodology}
% ============================================================================

We organize our analysis around three research questions, each employing complementary statistical methods.

\subsection{RQ1: Martingale Efficiency and Microstructure Noise}
\label{sec:method_efficiency}

\paragraph{Variance Ratio Tests.}
Following \citet{lo1988stock}, we compute the variance ratio $VR(q) = \text{Var}(p_t - p_{t-q}) / [q \cdot \text{Var}(p_t - p_{t-1})]$ for $q \in \{2, 5, 10, 20\}$ with heteroskedasticity-robust standard errors. Under the martingale hypothesis, $VR(q) = 1$.

\paragraph{Roll's Implied Spread.}
We estimate effective bid-ask spreads using \citet{roll1984simple}: $\hat{s} = 2\sqrt{-\text{Cov}(\Delta p_t, \Delta p_{t-1})}$. This recovers the spread from the negative first-order autocovariance induced by bid-ask bounce.

\paragraph{BDS Test.}
Following \citet{brock1996test}, we test for nonlinear dependence in AR(1) residuals of price changes. The BDS statistic tests whether the correlation integral at embedding dimension $m$ equals the $m$-th power of the correlation integral at dimension 1, as expected under i.i.d.

\subsection{RQ2: Extreme Value Characterization}
\label{sec:method_evt}

\paragraph{Generalized Pareto Distribution.}
For each market, we fit a GPD to exceedances above the 95th percentile of absolute returns using maximum likelihood \citep{coles2001introduction,mcneil2000estimation}:
\begin{equation}
    F_u(y) = 1 - \left(1 + \frac{\xi y}{\sigma}\right)^{-1/\xi}_+, \quad y > 0
\end{equation}
where $\xi$ is the shape parameter (tail index) and $\sigma$ the scale. Negative $\xi$ indicates bounded (Weibull-domain) tails; positive $\xi$ indicates heavy (Fr\'{e}chet-domain) tails.

\paragraph{Hill Estimator.}
We compute the non-parametric tail index $\hat{\alpha} = k / \sum_{i=1}^{k} \ln(X_{(i)}/X_{(k)})$ across order statistics $k \in [100, 10000]$ and report Hill plots for stability assessment.

\subsection{RQ3: Automated Tail Risk Strategy Optimization}
\label{sec:method_autoresearch}

To assess the economic significance of microstructure-based tail risk prediction, we adapt the automated hypothesis-testing framework of \citet{karpathy2025autoresearch} to systematically explore the model and feature space.

\paragraph{Framework.}
The optimization loop consists of three components:
\begin{enumerate}
    \item A \textbf{fixed backtester} that loads temporal hold-out data, defines tail events as $|\Delta p_t| > \text{VaR}_{99}$ (estimated from training data), executes the strategy with realistic transaction costs (Kalshi: 0.6\% per contract; Polymarket: 0.02\% + slippage), and computes the out-of-sample Sharpe ratio.
    \item An \textbf{editable strategy module} containing the predictive model (feature engineering, architecture, hyperparameters) and position sizing logic.
    \item An \textbf{experiment loop} that iterates: (a) form hypothesis, (b) modify strategy, (c) commit to version control, (d) evaluate, (e) keep if Sharpe improves, revert otherwise.
\end{enumerate}

\paragraph{Data Splits.}
We use a strict temporal split to prevent look-ahead bias:
\begin{itemize}
    \item \textbf{Train:} all trades before June 1, 2025
    \item \textbf{Validation:} June 1 -- August 31, 2025
    \item \textbf{Test:} September 1 -- November 25, 2025
\end{itemize}

\paragraph{Base Features.}
The fixed backtester provides \NBaseFeatures{} features per trade: current price, trade size, taker side, lagged returns ($\Delta p_{t-1}$, $\Delta p_{t-2}$), lagged trade size and taker side, 20-trade buy-side imbalance, 10- and 50-trade rolling volatility, hour of day (UTC), day of week, and distance to price boundary ($\min(p, 1-p)$).

\paragraph{Strategy Space.}
The agent explores: (a) derived features (interactions, polynomial terms, regime indicators); (b) model architectures (gradient boosting, logistic regression, ensembles, neural networks); (c) position sizing (proportional reduction, Kelly criterion, volatility targeting, stop-loss rules); and (d) advanced methods (GMM regime detection, copula features, reinforcement learning).

% ============================================================================
\section{Results}
\label{sec:results}
% ============================================================================

\subsection{Descriptive Statistics}

\input{tables/trade_summary.tex}
\input{tables/price_distribution.tex}
\input{tables/price_concentration.tex}
\input{tables/trade_size_distribution.tex}
\input{tables/taker_side_analysis.tex}

\begin{figure}[!htbp]\centering
\includegraphics[width=\linewidth]{figures/price_distribution.png}
\caption{Price (implied probability) distributions by platform.}
\label{fig:price_dist}
\end{figure}

\begin{figure}[!htbp]\centering
\includegraphics[width=\linewidth]{figures/volume_by_day.png}
\caption{Daily trading activity and notional volume by platform.}
\label{fig:volume_by_day}
\end{figure}

\begin{figure}[!htbp]\centering
\includegraphics[width=\linewidth]{figures/trade_size_distribution.png}
\caption{Trade size distributions by platform (log scale).}
\label{fig:trade_size}
\end{figure}

\begin{figure}[!htbp]\centering
\includegraphics[width=\linewidth]{figures/hourly_pattern.png}
\caption{Trading activity by hour of day (UTC).}
\label{fig:hourly}
\end{figure}

\subsection{RQ1: Martingale Efficiency}

% [Phase 2 results: VR(2) = 0.59, Roll's spread 0.68-1.72 cents,
%  political markets efficient after AR(1), sports markets retain nonlinearity]

\input{tables/variance_ratio.tex}
\input{tables/roll_spread.tex}
\input{tables/bds_results.tex}

\begin{figure}[!htbp]\centering
\includegraphics[width=\linewidth]{figures/variance_ratio_by_lag.png}
\caption{Variance ratio by lag across market types. Horizontal line at $VR = 1$ indicates martingale.}
\label{fig:vr_by_lag}
\end{figure}

\subsection{RQ2: Tail Risk Characterization}

% [Phase 2+3 results: Kalshi xi = -1.53, Polymarket xi = -0.19,
%  bounded tails consistent with binary contracts, tick-size compression]

\input{tables/gpd_results.tex}

\begin{figure}[!htbp]\centering
\includegraphics[width=\linewidth]{figures/hill_plot.png}
\caption{Hill estimator plots by platform. Stable regions indicate reliable tail index estimates.}
\label{fig:hill_plot}
\end{figure}

\begin{figure}[!htbp]\centering
\includegraphics[width=\linewidth]{figures/gpd_tail_fit.png}
\caption{GPD tail fit: predicted vs.\ empirical exceedance probabilities.}
\label{fig:gpd_fit}
\end{figure}

\subsection{RQ3: Tail Risk Prediction and Strategy Optimization}

% [Autoresearch results: N iterations, best Sharpe, AUC, feature importances,
%  ablation table, calibration, equity curve]

\input{tables/autoresearch_ablation.tex}

\begin{figure}[!htbp]\centering
\includegraphics[width=\linewidth]{figures/shap_summary.png}
\caption{SHAP feature importance for the best tail risk predictor. Taker side dominates.}
\label{fig:shap}
\end{figure}

\begin{figure}[!htbp]\centering
\includegraphics[width=\linewidth]{figures/calibration_curve.png}
\caption{Calibration curve: predicted vs.\ observed tail event frequency.}
\label{fig:calibration}
\end{figure}

\begin{figure}[!htbp]\centering
\includegraphics[width=\linewidth]{figures/cumulative_returns.png}
\caption{Cumulative returns of the tail-risk-adjusted strategy vs.\ buy-and-hold benchmark (test period, net of transaction costs).}
\label{fig:equity_curve}
\end{figure}

% ============================================================================
\section{Discussion}
\label{sec:discussion}
% ============================================================================

% [Key findings synthesis:
%  1. Bid-ask bounce dominates apparent inefficiency
%  2. Sports vs political markets are structurally different
%  3. Taker side as informed trading signal (link to PIN)
%  4. Discrete ticks compress tails on Kalshi
%  5. Economic significance: modest but real Sharpe after costs]

% ============================================================================
\section{Conclusion}
\label{sec:conclusion}
% ============================================================================

% [Summary of three RQs, contributions, limitations, future work]

\bibliographystyle{plainnat}
\bibliography{references}

\end{document}
